\section{Systematic power: complete proofs}
\label{sec:annex-p15-systematic-power}

Use the finite ground vocabulary and the range-measure
conventions of \cref{sec:p15-systematic-power}.
The proofs preserve the distinction between the original
finite minimal-sentence construction and the generalized
measure-based construction.

\begin{annexitem}{annex:p15.power-normal-forms}{thm:p15-power-normal-forms}{All truth assignments and both normal forms}
\begin{proof}
For an assignment \(u\), every conjunct of \(D_v\)
is true exactly when \(u=v\). By De Morgan's law,
\(M_v\equiv\neg D_v\), so \(M_v\) is false
exactly at \(v\). It follows immediately that
the disjunction of the \(D_v\)'s for which
\(S\) is true agrees with \(S\) at every assignment.
The conjunction of the \(M_v\)'s for which
\(S\) is false has precisely those assignments
excluded and likewise agrees with \(S\).
Ground freedom in the equality-free relational
language identifies this Boolean agreement
with first-order equivalence.

Every \(D_v\) has exactly one satisfying
assignment. Omitting it from the maximal
normal form changes truth at \(v\);
including one not in that form also
changes truth there. This proves uniqueness
of the disjunct set. The same argument,
using the unique excluded assignment
of \(M_v\), proves uniqueness of the
conjunct set.

Distinct \(D_v,D_w\) cannot both be
true at one assignment. Distinct
\(M_v,M_w\) cannot both be false
at one assignment, so their disjunction
is valid. If a singular \(H\)
follows from \(M_v\), it is true
at every assignment except possibly
\(v\). It is therefore valid or
has exactly the same truth set
as \(M_v\). If \(H\entails D_v\),
its truth set is a subset of
\(\{v\}\), hence either empty
or that singleton. This proves
both extremal-content claims.

There are \(2^n\) assignments,
and membership or nonmembership
in \([S]\) partitions them.
The content family contains one
distinct minimal sentence for
each nonmember, proving the
counting identity.
For \(n=0\), there is one empty
assignment, its state description
is \(\top\) and its minimal
sentence is \(\bot\).
The same truth-set arguments
give every stated empty-case
convention.
\end{proof}
\end{annexitem}

\begin{annexitem}{annex:p15.power-optimal-split}{thm:p15-power-optimal-split}{The original optimal partition is unique}
\begin{proof}
Let \(K'=K_1'\mathbin{\dot\cup}K_2'\)
be an admissible partition, and choose
\(M_v\in K_2'\).
If \(T\not\entails M_v\), some
first-order structure satisfies
\(T\) and falsifies \(M_v\).
Its ground assignment is \(v\).
It satisfies every other minimal
sentence, and hence every member
of \(K_1'\), since \(M_v\) does
not belong to that disjoint part.
This would be a model of
\(\{T\}\cup K_1'\) falsifying the
required consequence \(M_v\).
The contradiction proves
\(T\entails M_v\).
This argument applies to every
member of \(K_2'\) and to a
quantified \(T\) without assigning
it a truth value on incomplete
ground information.

Thus every admissible second
part is contained in
\(K_T=\{M\in K':T\entails M\}\).
Conversely \(K_T\) is itself
an admissible second part,
because all its members follow
from \(T\) alone.
The largest possible ratio
is therefore \(|K_T|/|K'|\).
A proper subset of the finite
\(K_T\) has smaller cardinality,
so the maximizing partition
has exactly this second part
and its uniquely determined
complementary first part.
If \(T\) is inconsistent,
it entails all minimal
sentences and the score
is one, as the same argument
requires.
\end{proof}
\end{annexitem}

\begin{annexitem}{annex:p15.power-uniform-formula}{thm:p15-power-uniform-formula}{Content intersection and the ground projection}
\begin{proof}
For singular \(T\), a minimal sentence
\(M_v\) follows from \(T\) exactly
when \(v\) falsifies \(T\).
Thus the numerator family consists
of those \(M_v\) whose assignment
falsifies both \(T\) and \(S\).
These are exactly the assignments
falsifying \(T\lor S\).
Dividing both numerator and
denominator counts by \(2^n\)
gives the displayed ratio
of uniform content measures.

Equivalent singular conjunctions
have the same truth set and
the same unique minimal normal
form. They therefore have the
same content family and score.
This proves invariance under
regrouping, duplication and
logically equivalent reformulation
of the data within the fixed
language.

For arbitrary \(T\), put
\(I_T=\{v:T\entails M_v\}\).
The conjunction \(T^\flat\)
of those minimal sentences
is false exactly at \(I_T\).
By the preceding optimal-split
theorem, the numerator is
indexed by \(I_T\cap(\Omega\setminus[S])\).
This is precisely the false
set of \(T^\flat\lor S\),
so the same ratio with
\(T^\flat\) computes the
original measure.

This construction records
which ground assignments
cannot extend to a model
of \(T\). It need not
determine whether \(T\)
is true in every structure
extending an allowed assignment.
The explicit pair in the
finite-procedure proof below
shows that difference.
\end{proof}
\end{annexitem}

\begin{annexitem}{annex:p15.power-source-examples}{thm:p15-power-source-examples}{The one-name example and all three source regroupings}
\begin{proof}
With atoms \(P(a),Q(a)\), the conjunction
\(S=P(a)\land Q(a)\) has the single
true assignment \(11\). Its false
assignments \(00,01,10\) give exactly
the three displayed minimal sentences.
The universal conditional rules out
the ground assignment \(10\), and
therefore entails \(\neg P(a)\lor Q(a)\).
It does not rule out \(00\) or \(01\):
each extends to a one-object model
of the universal conditional.
The numerator consequently has one
member and the denominator three.

For \(K_1\), the consequences
\(Q(a)\) and \(\neg P(b)\)
follow from \(P(a)\) and
\(\neg Q(b)\), respectively,
using the universal conditional.
Each other listed sentence can
be falsified while keeping \(T\)
and the remaining sentences true:
for \(P(a)\), choose the local
pair \((P(a),Q(a))=(0,1)\);
for \(\neg Q(b)\), choose
\((P(b),Q(b))=(0,1)\);
for \(Q(c)\), choose \((0,0)\)
at \(c\); and for \(\neg P(d)\),
choose \((1,1)\) at \(d\).
At all other names retain
the required literals and
choose any compatible pair.
Each pair satisfies \(P\to Q\),
so these assignments extend
to a model of \(T\).
Exactly two of six listed
sentences are thus individually
derivable from the remainder.

For \(K_2\), falsify its
first conjunction by choosing
the pair \((0,0)\) at \(a\).
Falsify its second conjunction
by choosing \((1,1)\) at \(b\).
Falsify its third and fourth
sentences with \((0,0)\) at
\(c\) and \((1,1)\) at \(d\),
respectively. Each choice
can retain the other three
sentences and \(T\).
Thus its naive fraction
is zero.

For \(K_3\), the second
conjunction follows from
the first using \(T\).
The first does not follow
from the others: use
\((0,1)\) at both \(a,b\),
keeping the second conjunction
true but the first false.
The \(c,d\) countermodels
just used apply to the
remaining two sentences.
Thus the naive fraction
is one fourth.

All three conjunctions
specify the same six
signed atoms. Of the
eight ground atoms
over four names, two
remain free, so exactly
four of \(256\) assignments
satisfy the data.
Its content family has
\(256-4=252\) members.
At each of the four
names, precisely three
of the four \(P,Q\)
assignments satisfy
the universal's instance.
Every resulting combination
extends to a four-object
model of \(T\), and any
violating instance cannot.
There are therefore
\(3^4=81\) compatible
ground assignments and
\(256-81=175\) minimal
sentences entailed by \(T\).
Every data-satisfying assignment
already satisfies those four
instances, so all \(175\)
entailed minimal sentences
belong to the data content.
The common score is
\(175/252=25/36\).
\end{proof}
\end{annexitem}

\begin{annexitem}{annex:p15.power-range-laws}{thm:p15-power-range-laws}{Deriving the measure laws from the stated axioms}
\begin{proof}
For any contradictory \(B\),
\(\top\lor B\) is valid
and \(\top\land B\)
is contradictory. Additivity
gives \(1=1+r(B)\),
so \(r(B)=0\).
In particular \(r(\bot)=0\).
Since \(U\lor\neg U\)
is valid and its
two disjuncts are
incompatible,
\[
                         r(U)+r(\neg U)=1 .
\]

Suppose \(U\equiv V\).
Both \(U\lor\neg U\)
and \(V\lor\neg U\)
are valid, with incompatible
disjuncts in each case.
Thus
\(r(U)+r(\neg U)=1
=r(V)+r(\neg U)\),
proving equivalence invariance
from the stated axioms.

If \(U\entails V\),
then \(V\equiv U\lor(V\land\neg U)\),
a disjoint decomposition.
Its second summand has
nonnegative measure, proving
\(r(U)\le r(V)\).
For arbitrary \(U,V\),
the disjoint decompositions
\[
 U\equiv(U\land V)\lor(U\land\neg V),
 \qquad
 U\lor V\equiv(U\land\neg V)\lor V
\]
give inclusion and exclusion
by eliminating \(r(U\land\neg V)\).
The definition \(g=1-r\)
and the complement formula
give \(g(U)=r(\neg U)\),
the interval bound and
\(g(\bot)=1\).
If \(U\lor V\) is
valid, then \(\neg U\)
and \(\neg V\) are
incompatible, so
\[
 g(U\land V)=r(\neg U\lor\neg V)
                           =g(U)+g(V).
\]
This proves every stated
content axiom.

For the examples, choose
structures \(\mathfrak M_1,\ldots,\mathfrak M_k\)
and nonnegative weights
\(w_j\) summing to one.
The function
\[
 r(U)=\sum_j w_j\mathbf1_{\{\mathfrak M_j\models U\}}
\]
is between zero and one,
equals one on valid
sentences and is additive
on incompatible disjunctions,
pointwise in each structure.
Likewise weights on the
finite ground assignments
give the axioms for
singular sentences.
Choosing one realizing
structure for each of
finitely many assignments
also gives a full-language
extension, although different
choices can give different
values to quantified sentences.
\end{proof}
\end{annexitem}

\begin{annexitem}{annex:p15.power-probability-dual}{thm:p15-power-probability-dual}{The exact dual conditional}
\begin{proof}
Complementation and De Morgan's law give
\[
 \frac{g(T\lor S)}{g(S)}
 =\frac{r(\neg(T\lor S))}{r(\neg S)}
 =\frac{r(\neg T\land\neg S)}{r(\neg S)}
 =p_r(\neg T,\neg S).
\]
The denominator is positive
precisely when \(g(S)>0\);
thus the identity never
uses a conditional outside
its domain.

For \(p_r(T,S)\), the
event or logical range
\(T\land S\) is contained
in \(S\), so monotonicity
gives a numerator between
zero and \(r(S)\).
Division by positive
\(r(S)\) gives the
unit-interval bound.
The identity just proved
gives the same bound
for \(s_r\).
Its two negated arguments
establish the stated
interpretation and distinguish
it from \(p_r(T,S)\).
\end{proof}
\end{annexitem}

\begin{annexitem}{annex:p15.power-eight-identities}{thm:p15-power-eight-identities}{Every identity and every denominator condition}
\begin{proof}
The two interval bounds
were proved above.
The disjoint decomposition
\[
 S\equiv(T\land S)\lor(\neg T\land S)
\]
gives the probability
complement identity after
division by \(r(S)>0\).
For the power complement,
use the duality formula:
\[
 s_r(\neg T,S)=p_r(T,\neg S)
              =1-p_r(\neg T,\neg S)
              =1-s_r(T,S).
\]
All terms here have
the same positive
denominator \(r(\neg S)\).

Apply range inclusion
and exclusion to
\(T_1\land S\) and
\(T_2\land S\).
Their union is
\((T_1\lor T_2)\land S\)
and intersection
\(T_1\land T_2\land S\).
Division by \(r(S)\)
proves the probability
addition formula.
Replace \(T_1,T_2,S\)
by their negations and
use the duality identity
and De Morgan's laws.
This proves the power
addition formula with
conjunction and disjunction
exchanged.

If \(r(S)>0\) and
\(r(T_1\land S)>0\),
multiply the defining ratios:
\[
 \frac{r(T_1\land S)}{r(S)}
 \frac{r(T_2\land T_1\land S)}{r(T_1\land S)}
 =\frac{r(T_1\land T_2\land S)}{r(S)}.
\]
This gives the probability
product formula.
The identical calculation
for the dual conditional,
or replacement of all
three sentences by
their negations, gives
the power product formula.
Its additional denominator
is \(r(\neg T_1\land\neg S)
=g(T_1\lor S)\).

If \(r(T_1\land S)=0\),
monotonicity forces
\(r(T_1\land T_2\land S)=0\).
The left conditional
is zero when \(r(S)>0\),
but the second right
conditional is not defined.
If \(g(T_1\lor S)=0\),
then
\(\neg(T_1\lor T_2\lor S)\)
has measure zero as
a subrange of
\(\neg(T_1\lor S)\).
The left power is zero
when \(g(S)>0\),
and the second right
power is undefined.
These cases establish
the required boundary
statements without a
convention for multiplying
an undefined quantity.
\end{proof}
\end{annexitem}

\begin{annexitem}{annex:p15.power-dualism}{thm:p15-power-dualism}{The full Boolean transformation}
\begin{proof}
For any Boolean expression \(B\)
in the sentence metavariables,
structural induction gives
\[
 D(B)(A_1,\ldots,A_m)
       \equiv\neg B(\neg A_1,\ldots,\neg A_m).
\]
For a metavariable this
is double negation.
For a negation it
follows from the
inductive hypothesis.
For a conjunction or
disjunction it is
De Morgan's law.
For a truth constant
it requires exactly
the stated exchange
of \(\top,\bot\).
These cases prove
the identity for all
Boolean arguments.

It follows that
\[
 s_r(D(B),D(C))
 =p_r(\neg D(B),\neg D(C))
 =p_r(B(\neg A_1,\ldots,\neg A_m),
      C(\neg A_1,\ldots,\neg A_m)).
\]
The power denominator
\(g(D(C))\) equals
\(r(C(\neg A_1,\ldots,\neg A_m))\),
so definedness is
preserved exactly.
Apply a universally
valid probability
equality or inequality
to the negated choices
of its sentence
metavariables. Replace
each resulting term
by the displayed
equal power term.
Numerical operations
and comparisons
are unchanged, proving
the desired assertion.

Complementation is
a bijection on the
sentence-equivalence
classes, and applying
\(D\) twice returns
the original Boolean
expression. The same
argument proves the
converse transformation.

For an additional
entailment premise,
the structural identity
changes it to
\(\neg D(B)\entails\neg D(C)\).
Classical contraposition
makes this equivalent
to \(D(C)\entails D(B)\),
with the direction
reversed. In particular,
leaving all logical
constants and hypotheses
unchanged would not
be a valid dualization:
\(p_r(\top,S)=1\)
does not become
\(s_r(\top,S)=1\),
which is false on
the latter's domain.
\end{proof}
\end{annexitem}

\begin{annexitem}{annex:p15.power-strength-boundary}{thm:p15-power-strength-boundary}{Strength does not certify a sound explanation}
\begin{proof}
Since \(g(\top)=0\),
\(\top\lor S\) has
zero content and
\(s_r(\top,S)=0\).
Since \(\bot\lor S\equiv S\),
\(s_r(\bot,S)=1\).
Likewise \(S\lor S\equiv S\)
gives \(s_r(S,S)=1\)
whenever the denominator
is positive.

If \(T\entails U\),
then \(\neg U\entails\neg T\).
Intersecting with \(\neg S\),
using measure monotonicity
and dividing by positive
\(r(\neg S)\) proves
\(s_r(T,S)\ge s_r(U,S)\).
If \(T\entails S\),
then \(\neg S\entails\neg T\),
so the dual conditional
\(p_r(\neg T,\neg S)\)
is one.

For finite singular
sentences with all
assignment weights
positive,
\[
 s_r(T,S)=1
 \quad\Longleftrightarrow\quad
 r(T\land\neg S)=0
 \quad\Longleftrightarrow\quad
 [T\land\neg S]=\varnothing.
\]
The last equivalence
uses positivity at
every formal assignment.
Ground freedom makes
the last condition
equivalent to \(T\entails S\).
The same argument
with \(S\land\neg T\)
proves the probability-one
counterpart.

For failure without
full support, take
two atoms \(A,B\)
and put probability
one half on \(00\),
one half on \(11\),
and zero on \(01,10\).
For \(T=A,S=B\)
both denominators
are one half and
both \(p_r(A,B)\)
and \(s_r(A,B)\)
equal one.
Yet the zero-weight
assignment \(01\)
refutes \(B\entails A\),
and \(10\) refutes
\(A\entails B\).
The maximum score
of an inconsistent
\(T\) and the
self-systematization
identity provide
additional reasons
not to identify
the score with
truth or independent
explanatory adequacy.
\end{proof}
\end{annexitem}

\begin{annexitem}{annex:p15.power-weight-choice}{thm:p15-power-weight-choice}{Uniform and latent-mixture weights}
\begin{proof}
With four equal weights,
\(r(A)=r(\neg A)=1/2\)
and
\(r(A\land B)=r(\neg A\land\neg B)=1/4\).
The defining conditional
ratios therefore both
equal one half.

For the latent construction,
conditional independence
gives probabilities
\(\theta^2\),
\(\theta(1-\theta)\),
\((1-\theta)\theta\)
and \((1-\theta)^2\)
for assignments \(11,10,01,00\), respectively.
Averaging the two
values of \(\theta\)
gives \(5/16\)
for \(00,11\)
and \(3/16\)
for \(01,10\).
The weights are
positive and sum
to one. Both
marginals still
equal one half.
Thus
\[
 p_r(B,A)=\frac{5/16}{1/2}=\frac58,
 \qquad
 s_r(B,A)=p_r(\neg B,\neg A)=\frac58.
\]
They define range
measures by the
proved finite-weight
construction. No
axiom among the
stated range axioms
chooses between
these two measures.
\end{proof}
\end{annexitem}

\begin{annexitem}{annex:p15.power-language-choice}{thm:p15-power-language-choice}{Vocabulary extensions and their limits}
\begin{proof}
If \(k\) fresh atoms
are added and neither
singular \(T\) nor
singular \(S\) uses
them, every old
assignment has
\(2^k\) extensions.
The numbers falsifying
\(S\) and falsifying
both \(T,S\) are
both multiplied
by \(2^k\).
Their ratio is
unchanged.

For \(T=A,S=B\)
on two uniform
atoms, exactly
one of the two
assignments with
\(B=0\) also
has \(A=0\),
giving one half.
After replacing
\(A\) by \(C\lor D\)
and using three
uniform atoms,
exactly one of
the four assignments
with \(B=0\)
has \(C=D=0\).
The ratio is
one fourth.
The translation
changes the weight
of the old
\(A=0\) class;
it is not an
equivalent reformulation
under one fixed
measure.

For the quantified
example, the one-name
calculation already
gave one third.
With two names
there are sixteen
ground assignments.
Four satisfy
\(P(a)\land Q(a)\),
so the denominator
content has twelve
minimal sentences.
Exactly nine assignments
satisfy both named
instances of \(P\to Q\);
each extends to
a two-object model
of the universal
theory. The other
seven assignments
are excluded by
that theory.
Of those seven,
one has \(a\)'s
pair equal to
\(11\) and \(b\)'s
pair equal to
\(10\); it satisfies
the data and so
does not belong
to the data content.
The other six
are excluded by
both theory and
data. Thus the
score is \(6/12=1/2\).
The new named
instances are
constrained by the
universal theory,
so the fresh-unconstrained-
atom argument does
not apply.
\end{proof}
\end{annexitem}

\begin{annexitem}{annex:p15.power-finite-procedure}{thm:p15-power-finite-procedure}{Exact calculation and the missing quantified truth value}
\begin{proof}
Finite recursive Boolean
evaluation determines
the truth of \(T,S\)
at every ground
assignment. The uniform
formula proves that
the ratio of the
two stated counts
is exactly the
systematic score,
provided the denominator
is positive.
For rational probability
weights, the corresponding
finite sums and their
positive-denominator ratio
are exact rational
operations.

For quantified \(T\),
the minimal sentence
\(M_v\) follows
from \(T\) exactly
when no structure
satisfies \(T\land D_v\),
because \(M_v\equiv\neg D_v\).
The optimal-split
theorem then proves
the stated numerator
calculation, with
at most \(2^n\)
model-existence queries.
It does not decide
any such arbitrary
first-order query
by finite Boolean
evaluation alone.

For an explicit
example, take the
one-name language
with \(P,Q\) and
the assignment
\(P(a)=Q(a)=1\).
The one-object
structure with both
predicates true
satisfies
\(T=\forall x(P(x)\to Q(x))\).
Add an unnamed
object at which
\(P\) is true
and \(Q\) false,
leaving both values
at \(a\) unchanged.
The enlarged structure
satisfies the same
ground assignment
and falsifies \(T\).
Thus that assignment
has extensions on
both sides of
the quantified
claim, as asserted.
\end{proof}
\end{annexitem}
